\section{AGI的演进路线图}

\subsection{模拟世界 $\rightarrow$ 探索世界 $\rightarrow$ 归纳世界}

蒋大昕提出AGI的三阶段演进路线：

\begin{knowledgebox}{三阶段框架}
\begin{enumerate}
\item \textbf{模拟世界}：GPT-4O代表——多模态融合，为物理世界建模
\item \textbf{探索世界}：FSD V12代表——端到端控制，从数字走向物理
\item \textbf{归纳世界}：O1代表——System 2慢思考，发现规律和知识
\end{enumerate}
过去18个月在三个阶段上都取得了标志性突破。
\end{knowledgebox}

\subsection{Scaling的历史视角}

杨植麟从AI七八十年的历史中提炼出一个核心结论：

\begin{importantbox}{AI历史上唯一有效的就是Scaling}
无论是参数量、数据量还是计算量——更大总是更好。O1不只是量变，而是质变：它开辟了新的Scaling维度（RL Scaling），打破了``数据用完了怎么办''的焦虑。更重要的是，Self-play产生的数据理论上没有上限。
\end{importantbox}

\subsection{本章小结}

AGI的演进可以从``模拟-探索-归纳''三阶段理解。Scaling仍是核心引擎，O1开辟了RL Scaling新维度，打破了数据墙的限制。


